\documentclass[12pt]{article}
\usepackage{geometry}                % See geometry.pdf to learn the layout options. There are lots.
\geometry{letterpaper}                   % ... or a4paper or a5paper or ... 
%\geometry{landscape}                % Activate for for rotated page geometry
\usepackage[parfill]{parskip}    % Activate to begin paragraphs with an empty line rather than an indent
\usepackage{daves,fancyhdr,natbib,graphicx,dcolumn,amsmath,lastpage,url}
\usepackage{amsmath,amssymb,epstopdf,longtable}
\usepackage[final]{pdfpages}
\DeclareGraphicsRule{.tif}{png}{.png}{`convert #1 `dirname #1`/`basename #1 .tif`.png}
\pagestyle{fancy}
\lhead{CE 5319 Machine Learning for Civil Engineers}
\rhead{SUMMER 2025}
\lfoot{ES9}
\cfoot{}
\rfoot{Page \thepage\ of \pageref{LastPage}}
\renewcommand\headrulewidth{0pt}



\begin{document}
\begin{center}
{\textbf{{ CE 5319 Machine Learning for Civil Engineers} \\ {Exercise Set 9} \\ {Support Vector Regression (SVR) and Rating Curves } }}
\end{center}

\subsection*{Purpose(s)}
This exercise, you will:
\begin{itemize}
    \item Fit and visualize Support Vector Regression (SVR) models using \texttt{scikit-learn}
    \item Compare SVR to linear regression on nonlinear hydrologic datasets
    \item Analyze the impact of hyperparameter $\epsilon$ on prediction smoothness and flexibility
    \item Relate machine learning tools to hydrologic modeling practice
\end{itemize}







\subsection*{Provided Starter Code (Do Not Edit)}
\begin{verbatim}
# Simulate a nonlinear rating curve (Stage - Discharge)
import numpy as np
import matplotlib.pyplot as plt
from sklearn.preprocessing import StandardScaler

# Seed and simulate
np.random.seed(0)
stage = np.linspace(1.0, 3.5, 30).reshape(-1, 1)
discharge = 30 * (stage.flatten() - 0.8)**2 + 
            np.random.normal(0, 5, size=stage.shape[0])

# Plot the raw data
plt.scatter(stage, discharge, color='gray')
plt.xlabel("Stage (ft)")
plt.ylabel("Discharge (cfs)")
plt.title("Simulated Stage-Discharge Data")
plt.grid(True)
plt.tight_layout()
plt.show()
\end{verbatim}

\clearpage

\section*{\small{Exercise(s)}}

\begin{enumerate}

\item SVR Model Fit and Prediction (40 pts)
\begin{enumerate}
    \item Fit an SVR model to the scaled data using:
    \begin{itemize}
        \item \texttt{kernel='rbf'}
        \item \texttt{C=100}
        \item \texttt{epsilon=0.1}
    \end{itemize}
    \item Invert the scaling to retrieve predicted discharge in original units.
    \item Plot observed vs predicted discharge values.
    \item Label the axes and add a legend to your plot.
\end{enumerate}

\item Comparison with Linear Regression (30 pts)
\begin{enumerate}
    \item Fit a linear regression model to the same scaled data.
    \item Plot both SVR and linear regression predictions against observed values.
    \item Compute and report the \textbf{RMSE} of both models.
    \item Interpret the result: Which model fits better? Why?
\end{enumerate}

\item Exploring $\epsilon$ Sensitivity (20 pts)
\begin{enumerate}
    \item Fit SVR models using three different \texttt{epsilon} values: \texttt{0.01}, \texttt{0.1}, and \texttt{0.5}.
    \item Plot the predictions from all three models on the same graph.
    \item Include observed data in the plot for reference.
    \item Discuss how changing $\epsilon$ affects the fit and flexibility of the model.
\end{enumerate}

\item Reflection and Application (10 pts)
Write a short paragraph (4--5 sentences) answering the following:
\begin{itemize}
    \item How might SVR be useful for estimating rating curves in real hydrologic applications?
    \item What limitations would you consider before using SVR operationally?
\end{itemize}







\end{enumerate}


\end{document}  

\section*{Assignment: SVR Applications in Hydrologic Modeling}

\subsection*{Course: CE 5319 – Machine Learning Applications in Civil Engineering}
\textbf{Topic:} Support Vector Regression (SVR) and Rating Curves \\
\textbf{Estimated Time:} 2–3 hours \\
\textbf{Deliverable:} Completed Jupyter Notebook with code, plots, and responses to discussion questions


\subsection*{Submission Notes}
\begin{itemize}
    \item Submit a clean \texttt{.ipynb} notebook with all cells executed and saved.
    \item Make sure each plot is clearly labeled and easy to interpret.
    \item Use comments in code to explain your steps.
\end{itemize}